\begin{ouverture}[molecules]{2}{Identifier la matière}{Constitution et transformations}{Observer · Modéliser · Expérimenter · Expliquer}
\objectif{Distinguer corps pur et mélange, espèce et échantillon.}
\objectif{Croiser tests chimiques, masse volumique et chromatographie.}
\objectif{Calculer une composition massique ou volumique.}
\prerequis{Proportionnalité, conversions et lecture de graphiques.}
\end{ouverture}
\essentiel
\essSection{1}{Corps purs et mélanges}
Un échantillon est une portion de matière. Une espèce chimique correspond à un ensemble d'entités identiques. Un corps pur ne contient qu'une espèce ; un mélange en contient plusieurs. Un mélange homogène présente une seule phase observable à l'échelle choisie. Un liquide transparent peut donc être un mélange.
\begin{exemple}L'eau salée limpide est homogène ; l'eau et l'huile non miscibles forment deux phases après décantation. L'air est un mélange gazeux : environ 78 \% de diazote, 21 \% de dioxygène et 1 \% d'autres gaz en volume.\end{exemple}
\essSection{2}{Masse volumique et composition}
\begin{definition}[Une propriété de l'échantillon]$\rho=m/V$. Avec $m$ en kilogrammes et $V$ en mètres cubes, $\rho$ s'exprime en \unit{kg.m^{-3}}. L'eau liquide vaut environ \qty{1,00}{g.mL^{-1}}, soit \qty{1,00e3}{kg.m^{-3}} ; l'air est de l'ordre de \qty{1}{kg.m^{-3}}. Les valeurs dépendent des conditions.\end{definition}
La fraction massique d'une espèce est $w=m_{\mathrm{espèce}}/m_{\mathrm{mélange}}$ ; le pourcentage vaut $100w$. Les fractions de toutes les espèces somment à 1. Les volumes de liquides mélangés ne sont pas toujours additifs.
\schema[corrige]{densite}
\essSection{3}{Identifier avec plusieurs indices}
Le sulfate de cuivre anhydre blanc devient bleu en présence d'eau. Le dioxyde de carbone trouble l'eau de chaux ; le dioxygène ravive une bûchette incandescente ; une très petite quantité de dihydrogène donne une détonation caractéristique au contact d'une flamme. Ces tests relèvent du laboratoire encadré.
\begin{savoirfaire}[Lire une CCM]Déposer les échantillons sur la ligne de dépôt, au-dessus de l'éluant ; fermer la cuve ; laisser monter l'éluant ; sortir la plaque et repérer le front. Plusieurs taches dans une piste révèlent plusieurs espèces détectées. Une tache à même hauteur qu'une référence, dans les mêmes conditions, est un indice d'identité.\end{savoirfaire}
\schema[corrige]{ccm-identification}
Une température de changement d'état ou une masse volumique proche d'une référence ne suffit pas toujours à identifier une espèce sans ambiguïté.
\essSection{4}{Expérimenter et modéliser}
\begin{experience}[Identifier un liquide]Peser un récipient vide puis le même récipient contenant un volume connu de liquide ; soustraire les masses. Comparer la masse volumique à une table à la même température. Compléter par un test ou une mesure de température de changement d'état. Porter blouse et lunettes ; ne pas goûter ni sentir directement les produits.\end{experience}
\exercices
\rubrique{Prendre confiance}
\begin{blocLivre}{ex-2-1}\begin{exercice}[1]{Trier des échantillons}
On dispose de cuivre pur, d’eau sucrée limpide et d’eau mêlée à de l’huile.
\begin{questions}
\item Classer ces trois échantillons.
\item La transparence garantit-elle la pureté ?
\end{questions}
\end{exercice}\end{blocLivre}
\begin{blocLivre}{ex-2-2}\begin{exercice}[1]{Mesurer une masse volumique}
Le récipient vide pèse 42,6 g ; rempli de 20,0 mL d’un liquide, il pèse 58,4 g.
\schema{densite}
\begin{questions}
\item Calculer la masse du liquide.
\item Calculer sa masse volumique.
\end{questions}
\end{exercice}\end{blocLivre}
\begin{blocLivre}{ex-2-3}\begin{exercice}[1]{Reconnaître un gaz}
Un gaz recueilli lors d’une réaction trouble l’eau de chaux.
\begin{questions}
\item Quelle espèce est mise en évidence ?
\item Quel test distinguerait le dioxygène ?
\end{questions}
\end{exercice}\end{blocLivre}
\begin{blocLivre}{ex-2-4}\begin{exercice}[1]{Air dans une pièce}
Une pièce contient \qty{50}{m^3} d’air. On adopte 21 \% de dioxygène en volume.
\begin{questions}
\item Estimer le volume de dioxygène.
\item Nommer le constituant majoritaire.
\end{questions}
\end{exercice}\end{blocLivre}
\rubrique{S’entraîner}
\begin{blocLivre}{ex-2-5}\begin{exercice}[2]{Composition d’un alliage}
Un alliage de 250 g contient 75,0 g de zinc et le reste de cuivre.
\begin{questions}
\item Calculer le pourcentage massique de zinc.
\item Déterminer la masse et le pourcentage de cuivre.
\end{questions}
\end{exercice}\end{blocLivre}
\begin{blocLivre}{ex-2-6}\begin{exercice}[2]{Une piste de chromatographie}
La plaque présente les références R et S et le mélange M. Toutes les pistes ont été développées ensemble.
\schema{ccm-identification}
\begin{questions}
\item Quelles espèces de référence sont compatibles avec M ?
\item M paraît-il pur dans cette analyse ?
\end{questions}
\end{exercice}\end{blocLivre}
\begin{blocLivre}{ex-2-7}\begin{exercice}[2]{Choisir un matériau}
Un solide occupe 12,0 cm$^3$ et pèse 32,4 g. Références : aluminium 2,70 ; fer 7,87 en \unit{g.cm^{-3}}.
\begin{questions}
\item Calculer sa masse volumique et proposer un matériau.
\item Quelle réserve faut-il conserver ?
\end{questions}
\end{exercice}\end{blocLivre}
\begin{blocLivre}{ex-2-8}\begin{exercice}[2]{Du volume à la masse}
Une bonbonne contient 40,0 L d’air de masse volumique \qty{1,20}{g.L^{-1}} dans les conditions de l’expérience.
\begin{questions}
\item Calculer la masse d’air.
\item Peut-on déduire la masse de dioxygène en multipliant par 21 \% ?
\end{questions}
\end{exercice}\end{blocLivre}
\rubrique{Raisonner et approfondir}
\begin{blocLivre}{ex-2-9}\begin{exercice}[3]{Un bilan de mélange}
On mélange 40 g d’une solution contenant 10 \% de sel en masse et 60 g d’eau pure. Aucune perte n’a lieu.
\begin{questions}
\item Calculer la masse de sel conservée.
\item Déterminer le pourcentage massique final.
\end{questions}
\end{exercice}\end{blocLivre}
\begin{blocLivre}{ex-2-10}\begin{exercice}[3]{Une pureté à discuter}
Un liquide bout vers 78 °C ; sa masse volumique vaut \qty{0,80}{g.mL^{-1}}. Une référence d’éthanol est voisine de ces valeurs.
\begin{questions}
\item Peut-on affirmer qu’il s’agit d’éthanol pur ?
\item Proposer une observation complémentaire.
\end{questions}
\end{exercice}\end{blocLivre}
\begin{blocLivre}{ex-2-11}\begin{exercice}[3]{Préparer une composition}
Un alliage doit contenir 12,0 \% d’étain en masse pour une masse totale de 2,50 kg.
\begin{questions}
\item Déterminer les masses d’étain et de l’autre métal.
\item Pourquoi 12 \% en volume donnerait-il un autre calcul ?
\end{questions}
\end{exercice}\end{blocLivre}
\begin{blocLivre}{ex-2-12}\begin{exercice}[3]{Un objet creux}
Un objet fermé en métal pèse 54 g. Son volume extérieur est 30 cm$^3$ ; le métal a pour masse volumique \qty{2,7}{g.cm^{-3}}.
\begin{questions}
\item Calculer le volume de métal.
\item En déduire le volume de la cavité, en négligeant la masse du gaz.
\end{questions}
\end{exercice}\end{blocLivre}

